\clearpage
\VerdierPhysicalPageAnchor{0175}%
\begingroup\small\setlength{\parskip}{0pt}\setlength{\emergencystretch}{1.5em}
\setlength{\abovedisplayskip}{4pt plus 2pt minus 2pt}%
\setlength{\belowdisplayskip}{4pt plus 2pt minus 2pt}%

\VerdierSemanticAnchor{ch3d-1-1-2-cont}
Il résulte immédiatement de la définition des triangles distingués dans
$K(\mathcal A)$ (chap. I, 3.3.1) et des propriétés du foncteur cohomologique
canonique $H\colon K(\mathcal A)\longrightarrow\mathcal A$ que les
catégories $K^{*,*'}(\mathcal A)\subset K(\mathcal A)$ sont des
sous-catégories triangulées (chap. II, 2.1.5) pleines (strictement si
$*={}$ ``vide'') et saturées (chap. II, 2.1.6) de $K(\mathcal A)$.

\VerdierSemanticAnchor{ch3d-1-1-3}
\medskip\noindent\textbf{1.1.3.}\quad
Un complexe $X$ de $\mathcal A$ est dit \emph{acyclique} si, pour tout entier
$i$, $H^i(X)=0$. On désigne par $\operatorname{Ac}^{*}(\mathcal A)$ la
sous-catégorie pleine de $K^{*}(\mathcal A)$ définie par les complexes
acycliques qui sont des objets de $K^{*}(\mathcal A)$. Un morphisme
$X\xrightarrow{u}Y$ de $K(\mathcal A)$ (resp.\ de
$\operatorname{comp}(\mathcal A)$) est appelé un \emph{quasi-isomorphisme}
si, pour tout entier $i$, le morphisme
$H^i(X)\xrightarrow{H^i(u)}H^i(Y)$ est un isomorphisme. On désigne par
$\operatorname{Qis}^{*}(\mathcal A)$ (ou par $\operatorname{Qis}^{*}$
lorsqu'aucune confusion n'en résulte) l'ensemble des quasi-isomorphismes de
$K^{*}(\mathcal A)$. La sous-catégorie
$\operatorname{Ac}^{*}(\mathcal A)$ de $K^{*}(\mathcal A)$ est une
sous-catégorie triangulée strictement pleine et saturée de
$K^{*}(\mathcal A)$. L'ensemble $\operatorname{Qis}^{*}(\mathcal A)$ est un
système multiplicatif saturé de morphismes de $K^{*}(\mathcal A)$
(chap. II, 2.1.4). Dans un triangle distingué $(X,Y,Z,u,v,w)$ de
$K^{*}(\mathcal A)$, le morphisme $u$ est un quasi-isomorphisme si et
seulement si le complexe $Z$ est acyclique. Autrement dit (chap. II,
2.1.15), $\operatorname{Qis}^{*}(\mathcal A)$ est le système multiplicatif
de $K^{*}(\mathcal A)$ associé à la sous-catégorie triangulée pleine
$\operatorname{Ac}^{*}(\mathcal A)$.

\VerdierSemanticAnchor{ch3d-1-1-4}
\medskip\noindent\textbf{1.1.4.}\quad
Soit $M$ un ensemble d'objets de $\mathcal A$. Un complexe de $\mathcal A$
est dit \emph{de type $M$} si tous ses composants sont des éléments de $M$.
Un morphisme $u\colon X\longrightarrow Y$ de $K^{*}(\mathcal A)$ ou de
$\operatorname{comp}^{*}(\mathcal A)$ est appelé une \emph{résolution à
droite} (resp. \emph{à gauche}) de $X$ (resp. $Y$) \emph{de type $M$} si
$u$ est un quasi-isomorphisme et si $Y$ (resp. $X$) est un complexe de type
$M$. Le plus souvent, par abus de langage et quand le contexte ne prête à
aucune confusion, on supprime la mention ``à droite'' (resp. ``à gauche'').
De même le plus souvent, par abus de langage et quand le contexte ne prête
à aucune confusion, on désigne par résolution à droite (resp. à gauche)
d'un objet $Y$, le but (resp.\ la source) d'une résolution à droite (resp.\ à
gauche) de $Y$. On désigne par
$\operatorname{comp}_{\mathcal A}^{*}(M)$, ou simplement
$\operatorname{comp}^{*}(M)$ (resp. $K_{\mathcal A}^{*}(M)$, ou simplement
$K^{*}(M)$), la sous-catégorie pleine de
$\operatorname{comp}^{*}(\mathcal A)$ (resp. $K^{*}(\mathcal A)$) définie
par les complexes de type $M$.

\VerdierSemanticAnchor{ch3d-1-1-5}
\medskip\noindent\textbf{1.1.5.}\quad
Lorsque $M$ est l'ensemble des objets projectifs (resp.\ injectifs) de la
catégorie $\mathcal A$, les complexes dont les composants sont projectifs
(resp.\ injectifs) sont appelés, conformément aux conventions ci-dessus, les
\emph{complexes de type projectif} (resp.\ \emph{injectif}). Le lecteur
prendra soin de ne pas confondre les complexes de type projectif (resp.\
injectif) avec les objets projectifs (resp.\ injectifs) de la catégorie
$\operatorname{comp}(\mathcal A)$. On constate aisément qu'un objet de
$\operatorname{comp}(\mathcal A)$ est projectif (resp.\ injectif) si et
seulement s'il est de type projectif (resp.\ injectif) et homotope à zéro.
\endgroup
